\section{Diferenciabilidade}


%%%%%%%%%%%%%%%%%%%%%%%%%%%%%%%%%%%%%%%%%%%%%%%%%%%%%%%%%%%%%%%%%%%
\subsection{Funções Diferenciáveis} 

A noção de função diferenciável, é devido a Fréchet e Stolz. 
Para $f:U\to\mathbb{R}$, com $U \subset \mathbb{R}^n$, e seja $a\in U$. A função $f$ é diferenciável no ponto $a$ quando existirem constantes $T_1, T_2,\ldots,T_n$, tais que para todo vetor $v = (\alpha_1,\ldots,\alpha_n) \in \mathbb{R}^n$, com $a+v\in U$, se tenha:
%
\[
f(a+v) = f(a) + \sum_{i=1}^{n}T_i\alpha_i + r(v) 
\ \ , \ \ 
\lim_{v\to 0} \frac{r(v)}{|v|} = 0
\]
%
Quando $f$ é diferenciável em todos os pontos de U, então $f$ é dita diferenciável. No caso particular em que $v = t e_i$, onde 
$(e_i)_{i=1}^{n}$ é a base canônica do $\mathbb{R}^n$, e retomamos 
a definição de derivada parcial, onde temos que $T_i$ é a derivada 
parcial de $f$ no ponto $a$:
%
\[
T_i = \frac{\partial f}{\partial x_i}(a) 
= \lim_{t \to 0} \frac{f(a + te_i) - f(a)}{t}
\]
%
Dizemos então que a função $f$ é diferenciável no ponto $a \in U$ 
quando existirem as derivadas parciais $T_i(a) = \partial_{i}f(a)$, 
para $i=1,\ldots,n$. E para todo vetor $v = (\alpha_1,\ldots,\alpha_n)$ tal que $a+v\in U$, tivermos:
%
\[
f(a+v) = f(a) + \sum_{i=1}^{n}\frac{\partial f}{\partial x_i}(a)
\alpha_i + r(v) 
\ \ , \ \ 
\lim_{v\to 0} \frac{r(v)}{|v|} = 0
\]
%
Onde $r(v)$ é o resto, e pode ser definido a partir da equação acima. 
%
\[
r(v) = f(a+v) - f(a) 
- \sum_{i=1}^{n}\frac{\partial f}{\partial x_i}(a)\alpha_i
\]
%
Essa definição é possível para qualquer função que possua derivadas parciais, sendo a principal condição a ser verificada quando for necessário provar que uma função é diferenciável. Devendo-se observar que o resultado $\lim_{v\to 0} \frac{r(v)}{|v|} = 0$ independe da norma utilizada. Mais a frente iremos discutir mais a fundo essa definição, mas já adiantando, o termo na somatória é nada mais do que o produto interno entre o gradiente da função $f$ no ponto $a$, e o vetor $v = (\alpha_1,\ldots,\alpha_n)$:
%
\[
\langle \nabla f(a) , v \rangle 
= \sum_{i=1}^{n}\frac{\partial f}{\partial x_i}(a)\alpha_i
\]
%
Esse resultado evidencia o motivo pelo qual enfatizamos a necessidade da dependência linear com a direção na derivada. Isso é resultado da linearidade do produto interno $\langle , \rangle$:
%
\[
\langle \nabla f(a) , tv + w \rangle 
= t \langle \nabla f(a) , v \rangle + \langle \nabla f(a) , w \rangle
\]



%%%%%%%%%%%%%%%%%%%%%%%%%%%%%%%%%%%%%%%%%%%%%%%%%%%%%%%%%%%%%%%%%%%
\subsection{Condição suficiente para diferenciabilidade} 

Em resumo, vimos que uma função é diferenciável no ponto se suas derivadas parciais existirem nesse ponto. Entretanto, a existência das derivadas parciais não implica necessariamente que $f$ seja uma função diferenciável no ponto. Um exemplo disso é a função:
%
\[
f(x,y) = \begin{cases} 
	\frac{xy^2}{x^2 + y^4} &, \text{se} \ x^2 + y^2 \neq 0 \\ 
	0                    &, \text{se} \  x = 0 \, y\neq 0 
\end{cases}
\]
%
Tomando o ponto $a$ como sendo a origem $(0,0)$, a função não é contínua nesse ponto, mas existem ambas derivadas parciais $\partial_x f(a)$ e $\partial_y f(a)$. 












